\chapter{Molecular Calculations}
\label{ch:molecular}

All molecular methods follow the same pattern: build the object from a
converged mean-field object of \pyscf{}, set attributes, call
\code{kernel()} or \code{run()}, and read the results from attributes.

\section{Choosing the method}

\begin{table}[!htb]
\centering
\small
\caption{Factory functions for molecules.}
\label{tab:factories}
\begin{tabularx}{\textwidth}{@{}l X l@{}}
\toprule
Call & Reference and class & Result \\
\midrule
\code{pycmf.OBMP2(mf)} & RHF or RKS $\to$ \code{ROBMP2}; UHF or ROHF $\to$ \code{UOBMP2} & OBMP2 \\
\code{pycmf.O2BMP2(mf, c_os=1.2)} & as above; \code{c_os} is mandatory & O2BMP2 \\
\code{pycmf.OBDH(mf, xc='B2PLYP')} & as above $\to$ \code{ROBDH} or \code{UOBDH} & one-body double hybrid \\
\code{pycmf.O2BDH(mf, c_os=1.2)} & as above; \code{c_os} is mandatory & opposite-spin OBDH \\
\code{pycmf.OBDF(mp, ncas, nelecas)} & converged OBMP2 or O2BMP2 object & active-space Hamiltonian \\
\bottomrule
\end{tabularx}
\end{table}

The factory functions choose the restricted or unrestricted class from the
type of the mean-field object. A ROHF reference is converted to UHF form.
Keyword arguments of the factory functions are set as attributes, so
\begin{lstlisting}[style=python]
mp = pycmf.OBMP2(mf, conv_tol=1e-10, level_shift=0.1)
\end{lstlisting}
is the same as setting \code{mp.conv_tol} and \code{mp.level_shift} after
construction. The classes \code{ROBMP2}, \code{UOBMP2}, \code{ROBDH} and
\code{UOBDH} can also be instantiated directly.

\section{OBMP2 and O2BMP2}
\label{sec:obmp2}

\subsection{A first calculation}

\begin{lstlisting}[style=python]
from pyscf import gto, scf, mp
import pycmf

mol = gto.M(atom='O 0 0 0.1173; H 0 0.7572 -0.4692; H 0 -0.7572 -0.4692',
            basis='cc-pvdz')
mf = scf.RHF(mol).run()
ob = pycmf.OBMP2(mf).run()
print(ob.e_tot, ob.converged)
print(mp.MP2(mf).run().e_tot)          # conventional MP2 for comparison
print(ob.ip_koopmans() * 27.2114)       # Koopmans IP in eV
\end{lstlisting}

For water in the cc-pVDZ basis this gives (from
\code{examples/01-obmp2.py}):

\begin{lstlisting}[style=output]
E(RHF)    = -76.0267720534
E(MP2)    = -76.2307756326
E(OBMP2)  = -76.2312896416   converged: True
Koopmans IP (eV) = 14.953
\end{lstlisting}

\subsection{Attributes}

\begin{longtable}{@{}p{0.27\textwidth} p{0.16\textwidth} p{0.49\textwidth}@{}}
\caption{Attributes of the molecular OBMP2 classes.}\label{tab:attrs}\\
\toprule
Attribute & Default & Meaning \\
\midrule
\endfirsthead
\toprule
Attribute & Default & Meaning \\
\midrule
\endhead
\bottomrule
\endfoot
\code{conv_tol} & \code{1e-8} & convergence threshold of the energy change between cycles \\
\code{conv_tol_grad} & $\sqrt{\text{\code{conv_tol}}}$ & threshold on the norm of the occupied--virtual block $\bar f_{ia}$ \\
\code{max_cycle} & \code{100} & maximum number of cycles; \code{max_cycle=1} from HF orbitals gives conventional MP2 \\
\code{bch_order} & \code{2} & 2: second-order BCH; 1: half-weighted first commutator \\
\code{c_os}, \code{c_ss} & \code{1, 1} & scaling of the opposite-spin and same-spin amplitudes \\
\code{level_shift} & \code{0} & constant added to the energy gap in the amplitude denominators \\
\code{occupation} & \code{'aufbau'} & \code{'mom'} or \code{'imom'} for state-specific calculations \\
\code{diis} & \code{True} & DIIS extrapolation of the correlated Fock matrix \\
\code{diis_space} & \code{8} & number of stored DIIS vectors \\
\code{diis_start_cycle} & \code{1} & first cycle with DIIS \\
\code{with_df} & from \code{mf} & density-fitting object; \code{None} means exact integrals \\
\code{max_memory} & from \code{mf} & memory limit in MB \\
\code{verbose}, \code{stdout} & from \code{mol} & output level and stream (\pyscf{} logger) \\
\end{longtable}

\code{frozen} orbitals are not supported for molecules: all orbitals are
relaxed by the correlated Fock operator, and a value other than \code{None}
or \code{0} raises \code{NotImplementedError}.

\subsection{Results}

After \code{kernel()} the following attributes are set.

\begin{description}[style=nextline,leftmargin=3.2cm,labelwidth=3cm]
  \item[\code{converged}] \code{True} if both convergence criteria were met.
  \item[\code{e_tot}] total energy; also the return value of \code{kernel()}.
  \item[\code{e_corr}] \code{e_tot} minus the energy of the input mean field.
  \item[\code{e_ref}] HF (or, for OBDH, Kohn--Sham) energy of the final determinant.
  \item[\code{e_hf}] energy of the input mean field.
  \item[\code{mo_energy}] eigenvalues of the correlated Fock matrix.
  \item[\code{mo_coeff}] correlated orbitals, occupied orbitals first.
  \item[\code{mo_occ}] occupation numbers in the same order.
\end{description}

For an unrestricted reference these arrays have a leading spin index of
length two, as in \pyscf{}.

\subsection{Methods of the result object}

\begin{longtable}{@{}p{0.40\textwidth} p{0.54\textwidth}@{}}
\toprule
Method & Returns \\
\midrule
\endhead
\bottomrule
\endfoot
\code{get_fock(ao_repr=False)} & converged correlated Fock matrix (MO or AO basis) \\
\code{get_vcorr()} & correlation potential $v$ in the basis of the converged orbitals \\
\code{make_rdm1(ao_repr=True, correlated=False)} & one-particle density; \code{correlated=True} gives the MP2-like density \\
\code{dip_moment(correlated=False, unit='Debye')} & dipole moment over either density \\
\code{ip_koopmans()}, \code{ea_koopmans()} & $-\varepsilon_{\mathrm{HOMO}}$ and $-\varepsilon_{\mathrm{LUMO}}$ of $\fbar$ (per spin for unrestricted) \\
\code{spin_square(correlated=False)} & $\langle S^2\rangle$ and $2S+1$ (unrestricted only; \code{ROBMP2} returns \code{(0, 1)}) \\
\code{hfcc(atoms=None, correlated=False)} & isotropic hyperfine couplings in MHz (unrestricted only) \\
\code{get_intermediates(with_amplitudes=False)} & quantities of the last cycle (\code{Intermediates}); amplitudes are rebuilt on request \\
\code{density_fit(auxbasis=None, with_df=None)} & switches the correlation part to density fitting (returns \code{self}) \\
\end{longtable}

\subsection{O2BMP2 and spin-component scaling}

\begin{lstlisting}[style=python]
o2 = pycmf.O2BMP2(mf, c_os=1.2).run()            # c_ss = 0
scs = pycmf.OBMP2(mf, c_os=1.2, c_ss=1/3).run()  # any other scaling
\end{lstlisting}

\code{O2BMP2} raises \code{ValueError} if \code{c_os} is not given. With
density fitting an O2BMP2 cycle costs about 60--70\,\% of an OBMP2 cycle
because the same-spin blocks are skipped. The example
\code{examples/02-o2bmp2.py} shows the effect of \code{c_os} on the
singlet--triplet gap of CH$_2$ (cc-pVTZ, density fitting):

\begin{lstlisting}[style=output]
method             E(S)-E(T) kcal/mol
HF                      28.24
OBMP2                   16.16
O2BMP2 c_os=1.1         12.09
O2BMP2 c_os=1.2         12.35
O2BMP2 c_os=1.3         12.90
\end{lstlisting}

\subsection{Open shells}

An unrestricted reference gives unrestricted OBMP2. OBMP2 reduces the spin
contamination of UHF considerably. For the CN radical in the cc-pVDZ basis
(\code{examples/06-properties.py}), $\langle S^2\rangle$ drops from 1.150
(UHF) to 0.766 (OBMP2 determinant) and 0.761 with the amplitude correction.

\begin{lstlisting}[style=python]
mol = gto.M(atom='O 0 0 0; H 0 0 0.97', basis='cc-pvdz', spin=1)
mf = scf.UHF(mol).run()
ob = pycmf.OBMP2(mf).run()
print(ob.e_tot, ob.spin_square())
\end{lstlisting}

\section{Density fitting}
\label{sec:df}

There are two ways to use density-fitted integrals.

\begin{enumerate}
  \item The mean field uses density fitting:
        \code{mf = scf.RHF(mol).density_fit()}. \pycmf{} reuses
        \code{mf.with_df} (JK-fit basis) for the correlation part, as the
        DF-MP2 of \pyscf{} does.
  \item The mean field uses exact integrals:
        \code{pycmf.OBMP2(mf).density_fit()} fits only the correlation part,
        with an RI-C auxiliary basis chosen by \pyscf{};
        \code{density_fit(auxbasis='cc-pvtz-ri')} selects a specific one.
\end{enumerate}

\begin{lstlisting}[style=python]
mf_df = scf.RHF(mol).density_fit().run()
e1 = pycmf.OBMP2(mf_df).kernel()                              # JK-fit basis
e2 = pycmf.OBMP2(mf).density_fit(auxbasis='cc-pvdz-ri').kernel()   # RI-C basis
\end{lstlisting}

For water/cc-pVDZ the two give $-76.2312271301$ and $-76.2312761456$~Hartree,
against $-76.2312896416$ with exact integrals. Exact integrals need the AO
integrals in memory; for more than a few hundred basis functions use density
fitting.

\section{Convergence}
\label{sec:convergence}

Every method uses the same self-consistent driver
(\code{pycmf.core.scf_loop}). At each cycle it

\begin{enumerate}
  \item builds the correlated Fock matrix from the current orbitals and
        orbital energies;
  \item extrapolates it with DIIS, using the commutator $[F,D]$ per spin
        channel as error vector;
  \item diagonalises it in the space spanned by the current orbitals;
  \item occupies the new orbitals according to \code{occupation}.
\end{enumerate}

Convergence requires both $|\Delta E| <$ \code{conv_tol} and
$\|\bar f_{ia}\| <$ \code{conv_tol_grad}. At \code{verbose >= 4} the
settings are printed first, then the energy, $\Delta E$ and
$\|\bar f_{ia}\|$ of every cycle. For water/cc-pVDZ with the default settings:

\begin{lstlisting}[style=output]
******** <class 'pycmf.mol.obmp2.ROBMP2'> ********
bch_order = 2
c_os = 1  c_ss = 1
level_shift = 0
occupation = aufbau
conv_tol = 1e-08
max_cycle = 100
DIIS = True, space = 8
integrals = exact
cycle= 1 E= -76.2307756295253  delta_E= -76.2  |f_ov|= 0.0302
cycle= 2 E= -76.231238968644  delta_E= -0.000463  |f_ov|= 0.00485
cycle= 3 E= -76.2312803148426  delta_E= -4.13e-05  |f_ov|= 0.00109
cycle= 4 E= -76.2312879484375  delta_E= -7.63e-06  |f_ov|= 0.000263
cycle= 5 E= -76.2312891373978  delta_E= -1.19e-06  |f_ov|= 9.69e-05
cycle= 6 E= -76.231289619696  delta_E= -4.82e-07  |f_ov|= 1.47e-05
cycle= 7 E= -76.2312896408706  delta_E= -2.12e-08  |f_ov|= 1.6e-06
cycle= 8 E= -76.2312896416401  delta_E= -7.69e-10  |f_ov|= 1.93e-07
ROBMP2 converged.  E(ROBMP2) = -76.2312896416401
\end{lstlisting}

The first cycle reproduces the MP2 energy ($-76.23077563$) to the accuracy
of the HF reference.

If a calculation does not converge:

\begin{itemize}
  \item increase \code{max_cycle};
  \item try a small \code{level_shift} (for example 0.1--0.3~Hartree), which
        damps the amplitudes of near-degenerate pairs;
  \item start DIIS later (\code{diis_start_cycle}) or switch it off to see
        whether the iterations oscillate;
  \item check that the mean field is a stable solution
        (\code{mf.stability()} in \pyscf{}).
\end{itemize}

A calculation that stops at \code{max_cycle} prints a warning and sets
\code{converged=False}; its energy is that of the last cycle.

\section{One-body double hybrids}
\label{sec:obdh}

\code{OBDH} and \code{O2BDH} take a converged HF or Kohn--Sham object, which
provides the starting orbitals. The functional is defined in one of three
ways:

\begin{lstlisting}[style=python]
pycmf.OBDH(mf)                                           # B2PLYP
pycmf.OBDH(mf, xc='B2GP-PLYP')                           # by name
pycmf.OBDH(mf, ax=0.6, ac=0.3, xc_x='PBE', xc_c='PBE')   # four parameters
pycmf.OBDH(mf, xc='0.5*HF + 0.5*B88, 0.75*LYP', ac=0.25) # PySCF string + a_c
pycmf.O2BDH(mf, c_os=1.2, xc='PBE0-DH')
\end{lstlisting}

Names are case-insensitive, and hyphens and underscores are ignored. With a
functional string the semilocal correlation must already be scaled by
$1-a_c$, and \code{ax} must not be given (the HF fraction is part of the
string). Explicit \code{ax} and \code{ac} override the parameters of a named
functional.

\begin{longtable}{@{}p{0.22\textwidth} p{0.18\textwidth} p{0.52\textwidth}@{}}
\caption{Additional attributes of \code{ROBDH} and \code{UOBDH}.}\\
\toprule
Attribute & Default & Meaning \\
\midrule
\endhead
\bottomrule
\endfoot
\code{ax} & \code{0.53} & fraction of HF exchange \\
\code{ac} & \code{0.27} & weight of the OBMP2 correlation \\
\code{xc_x}, \code{xc_c} & \code{'B88'}, \code{'LYP'} & semilocal exchange and correlation \\
\code{xc} & \code{None} & name of a double hybrid or a \pyscf{} functional string \\
\code{grids_level} & \code{None} & DFT integration grid level (default of \pyscf{}) \\
\end{longtable}

Useful methods: \code{get_xc()} returns the functional string of the
Kohn--Sham part, \code{set_functional(name)} takes the four parameters from a
named double hybrid, and \code{get_ks()} returns the Kohn--Sham object that
evaluates the hybrid part. When the reference uses density fitting, the
Kohn--Sham part uses the same density-fitting object.

From \code{examples/05-obdh.py} (water, cc-pVDZ, density fitting):

\begin{lstlisting}[style=output]
functional          : 0.53*HF + 0.47*B88, 0.73*LYP  + 0.27 E_c(OBMP2)
E(OBDH)             = -76.3425040639   converged: True
E(O2BDH, c_os=1.2)  = -76.3231558818
B2GP-PLYP  a_x = 0.6500  a_c = 0.3600  E = -76.3273800422
PBE0-DH    a_x = 0.5000  a_c = 0.1250  E = -76.3282840566
PBE-QIDH   a_x = 0.6934  a_c = 0.3333  E = -76.3072423903
\end{lstlisting}

For degenerate open shells such as the $^2\Pi$ state of OH, the orientation
of the open-shell orbital is fixed only by the numerical noise of the DFT
grid. Convergence then stalls at $\|\bar f_{ia}\|\sim10^{-6}$ and total
energies are reproducible to about $10^{-7}$~Hartree only.

\section{State-specific excited states}
\label{sec:mom}

An excited state is computed from a non-Aufbau determinant. First converge
the excited determinant at the HF level with the MOM of \pyscf{}, then run
OBMP2 or O2BMP2 with \code{occupation='imom'}:

\begin{lstlisting}[style=python]
mol = gto.M(atom='O 0 0 0.1173; H 0 0.7572 -0.4692; H 0 -0.7572 -0.4692',
            basis='aug-cc-pvdz')
mf_gs = scf.UHF(mol).run()

occ = mf_gs.mo_occ.copy()                   # alpha HOMO -> LUMO
homo = mol.nelec[0] - 1
occ[0][homo], occ[0][homo + 1] = 0, 1
mf_ex = scf.addons.mom_occ(scf.UHF(mol), mf_gs.mo_coeff, occ)
mf_ex.kernel(mf_ex.make_rdm1(mf_gs.mo_coeff, occ))

gs = pycmf.O2BMP2(mf_gs, c_os=1.2).run()
ex = pycmf.O2BMP2(mf_ex, c_os=1.2, occupation='imom').run()
omega = (ex.e_tot - gs.e_tot) * 27.2114     # eV
\end{lstlisting}

The occupation numbers \code{mo_occ} of the mean-field object define the
starting determinant; the orbitals are reordered occupied-first internally.
At \code{verbose >= 5} the number of electrons that have left the reference
occupied space is printed every cycle, which reveals a collapse to another
state. \code{'imom'} keeps the reference fixed and is usually more robust
than \code{'mom'}. For the example above (\code{examples/03-excited_state.py}):

\begin{lstlisting}[style=output]
method             omega (eV)    <S^2>
MOM-HF                  6.209    1.008
OBMP2                   7.588    1.006   converged: True
O2BMP2 (1.2)            7.193    1.006   converged: True
\end{lstlisting}

The excited determinant is a 50:50 mixture of singlet and triplet. Spin
purification is left to the user.

\section{One-body downfolding}
\label{sec:obdf}

\code{OBDF} builds the active-space Hamiltonian of
Section~\ref{sec:theory-obdf} from a converged OBMP2 or O2BMP2 object (OBDH
objects are not accepted).

\begin{lstlisting}[style=python]
mp = pycmf.OBMP2(mf).run()
ham = pycmf.OBDF(mp, ncas=6, nelecas=6)        # orbitals around the Fermi level
ham = pycmf.OBDF(mp, 4, (1, 1), cas_list=[0, 1, 4, 5])   # 0-based indices
e = ham.kernel()                  # FCI in the active space
h1, e_core = ham.get_h1eff()      # dressed one-body integrals and constant
h2 = ham.get_h2eff()              # bare active-space integrals
ham.to_fcidump('active.fcidump')  # restricted reference only
\end{lstlisting}

\begin{description}[style=nextline,leftmargin=3.2cm,labelwidth=3cm]
  \item[\code{ncas}] number of active orbitals.
  \item[\code{nelecas}] number of active electrons, an integer or a pair
        $(n_\alpha, n_\beta)$.
  \item[\code{cas_list}] 0-based indices of the active orbitals in the orbital
        order of \code{mp}; a pair of lists for different $\alpha$ and
        $\beta$ choices. By default the \code{ncas} orbitals around the Fermi
        level that hold \code{nelecas} electrons are used.
  \item[\code{fcisolver}] solver with the interface of \code{pyscf.fci};
        it can be replaced by any compatible solver (DMRG, selected CI).
  \item[\code{nroots}] number of states requested from the solver.
\end{description}

Results: \code{e_core}, \code{h1eff}, \code{h2eff}, and after
\code{kernel()} the total energy \code{e_tot}, the active-space energy
\code{e_cas} and the CI vector(s) \code{ci}. For an unrestricted reference
\code{h1eff} is a pair and \code{h2eff} a triple ($\alpha\alpha$,
$\alpha\beta$, $\beta\beta$), and \code{to_fcidump} is not available.

Do not cut through a set of degenerate orbitals with \code{cas_list}: the
result then depends on an arbitrary rotation. \code{OBDF} warns in this
situation.

\code{examples/04-obdf.py} compares OBDF(6,6) with OBMP2, CASCI(6,6) and
CCSD(T) along the dissociation curve of N$_2$ (cc-pVDZ):

\begin{lstlisting}[style=output]
 R (A)          OBMP2     CASCI(6,6)      OBDF(6,6)        CCSD(T)
  1.00    -109.215012    -108.980201    -109.220021    -109.231138
  1.10    -109.267377    -109.021905    -109.272409    -109.279398
  1.30    -109.229093    -108.963955    -109.225534    -109.224738
  1.60    -109.122703    -108.828238    -109.073809    -109.076363
  2.00    -109.053390    -108.741134    -108.924800    -109.015388
\end{lstlisting}

\section{Properties}
\label{sec:properties}

\begin{lstlisting}[style=python]
mp.make_rdm1(ao_repr=True, correlated=False)  # HF-like; True: MP2-like
mp.dip_moment(correlated=True)
mp.spin_square(correlated=False)              # unrestricted only
mp.hfcc(atoms=None, correlated=False)         # isotropic couplings, MHz
mp.ip_koopmans(); mp.ea_koopmans()
mp.get_fock(); mp.get_vcorr()                 # MO basis
\end{lstlisting}

The functions behind these methods are in \code{pycmf.properties}
(\code{rdm}, \code{spin}, \code{hfcc}). From
\code{examples/06-properties.py} (CN, cc-pVDZ):

\begin{lstlisting}[style=output]
<S^2>: UHF 1.1497 | OBMP2 determinant 0.7657 | with amplitudes 0.7609
HF-like  density: dipole 1.2327 D   A_iso(13C)   547.33 MHz   A_iso(14N)   -18.43 MHz
MP2-like density: dipole 1.0926 D   A_iso(13C)   560.71 MHz   A_iso(14N)   -10.45 MHz
\end{lstlisting}

\section{Experimental: restricted open-shell OBMP2}
\label{sec:roobmp2}

\code{ROOBMP2} works with a high-spin ROHF reference and keeps common spatial
orbitals, ordered closed, open, virtual. It has no formal derivation yet, no
DIIS and no state following.

\begin{lstlisting}[style=python]
from pycmf.mol.roobmp2 import ROOBMP2
mf = scf.ROHF(mol).run()
ro = ROOBMP2(mf).run()
print(ro.e_tot, ro.spin_square())      # <S^2> = 0.75 for a doublet
\end{lstlisting}

The additional attribute \code{conv_tol_denom} (default $10^{-6}$) is the
threshold on the difference between the amplitude denominators and the
diagonals of the spin Fock matrices. \code{mo_energy} carries the spin
energies as the tags \code{mo_ea} and \code{mo_eb}.

\section{Experimental: embedding in DFT}
\label{sec:embedding}

\begin{lstlisting}[style=python]
from pyscf import dft
from pycmf.embed import EmbeddedOBDH, EmbeddedOBMP2

ks = dft.UKS(mol, xc='pbe').density_fit().run()   # low level, full system
emb = EmbeddedOBDH(ks, active_atoms=[0, 1, 2], xc='B2PLYP',
                   use_cl=True, n_shells=1)
e = emb.kernel()
emb.e_env, emb.e_sub, emb.nelec_sub
emb.make_rdm1(); emb.dip_moment()                 # full system
EmbeddedOBMP2(ks, active_atoms=[0, 1, 2]).kernel()
\end{lstlisting}

\begin{longtable}{@{}p{0.22\textwidth} p{0.18\textwidth} p{0.52\textwidth}@{}}
\caption{Attributes of \code{EmbeddedOBMP2} and \code{EmbeddedOBDH}.}\\
\toprule
Attribute & Default & Meaning \\
\midrule
\endhead
\bottomrule
\endfoot
\code{active_atoms} & \code{None} & indices of the atoms of the subsystem \\
\code{xc_env} & from \code{mf} & low-level functional (environment and embedding potential); \code{'pbe'} for an HF object \\
\code{mu} & \code{1e6} & level-shift parameter of the projector \\
\code{use_cl} & \code{False} & truncate the virtual space by concentric localization \\
\code{n_shells} & \code{1} & number of CL shells kept \\
\end{longtable}

Results: \code{e_tot} (embedded total energy), \code{e_env} (low-level
energy of the full system), \code{e_sub} (high-level energy of the
subsystem), \code{e_low_sub}, \code{nelec_sub}, and the orbitals of the
embedded subsystem. \code{make_rdm1(with_env=True)} returns the density of
the full system. The embedding is formulated for unrestricted references
only. The helper functions \code{spade_partition}, \code{atom_aos},
\code{concentric_localization} and \code{shell_sizes} are exported by
\code{pycmf.embed}.

In version 1.0 the two levels have separate functionals: the environment and
the embedding potential use \code{xc_env}, and the subsystem uses its own
double hybrid. In version 0.2 \code{xc_env} was also used for the
Kohn--Sham part of the subsystem; those numbers are reproduced with
\code{xc=xc_env}.
